\chapter{Resultados e discussão}\label{cap:resultados}

% Manual: não usar "Desenvolvimento" como título. Apresentar o software com prints,
% tabelas com título acima e fonte abaixo, toda figura citada no texto. Comparar
% com o Estado da Arte e ser honesto sobre as dificuldades.
% Todo número deste capítulo sai de execução real: até lá, usar \resultado.

\section{Desempenho do modelo}

\pendente{Tabela com acurácia, precisão, recall e F1 por classe no conjunto de teste próprio (\resultado), matriz de confusão e discussão das classes que mais se confundem.}

\section{O aplicativo}

\pendente{Uma subseção por tela (câmera, resultado, histórico), com print e descrição de cada função.}

\section{Testes funcionais e de desempenho}

\pendente{Resultado dos casos de teste, tempo de resposta por aparelho (\resultado) e funcionamento em modo avião.}

\section{Comparação com os trabalhos correlatos}

\pendente{Retomar o quadro do Capítulo~\ref{cap:estado-da-arte}: onde este trabalho avança, confirma ou diverge.}

\section{Dificuldades e limitações}

\pendente{Diferença entre as fotos do dataset e as de campo, subjetividade da classe ``de vez'', tamanho do conjunto de teste, entre outras.}
